% =====================================================================
% Section 9 body fragment: Exploratory High-Mach Sounding-Rocket Results
% Parent jsr_paper.tex declares the \section; this file is \input inside it.
% Numbers: rocket-flight-database/sounding_rockets_exploratory.csv (full set),
% framing per _DRAFTING_BRIEF.md sec. 0 (report all, no cherry-pick).
% =====================================================================

The 25-flight headline corpus of Sec.~8 is bounded above by Mach~4.33, so it
contains no hypersonic flights. To probe whether the pre-pass and its downstream
supersonic models continue to function beyond that ceiling, we assembled a
separate set of roughly twenty historical sounding-rocket flights with
documented apogees, several of which peak between Mach~5 and Mach~7. This set is
\emph{not} part of the validation corpus and is reported here only as an
exploratory probe of method reach. We state at the outset that these flights are
\emph{excluded from every headline statistic}: presenting only the three that
happen to fall within $\pm10\%$ would constitute outcome-based selection bias,
the single most serious threat to the credibility of the result. We therefore
report the full set---successes and failures alike---in Table~\ref{tab:exploratory}.

\subsection{Full Exploratory Set}

Of the historical flights simulated, three fall within $\pm10\%$ of the measured
apogee. The Black Brant~V~VB (AAF-VB-32, Churchill Research Range, 3~March~1971;
\cite{dtic_ad0733141}) reaches a peak Mach number of $7.22$ and is predicted to
within $-6.97\%$ of its $273.6\,\mathrm{km}$ measured apogee. The two
Nike--Deacon (DAN) meteorological flights of \cite{heitkotter1956}, peaking at
Mach~$4.96$ and Mach~$5.08$, are predicted to $-1.06\%$ and $-0.89\%$
respectively. These three results establish that the integrated method
\emph{reaches} Mach~7 and returns apogee predictions within about $7\%$
\emph{when the vehicle is well characterized}---that is, when an experimentally
documented thrust history and a complete external geometry are available, as they
are for these flights.

The remaining flights fall well outside $\pm10\%$, and they fail in two opposite
directions. The 1965 Nike--Apache set, whose vehicle definition follows the
NASA performance handbook \cite{nasa_x721_66_568} and whose measured apogees are
taken from the NASA Wallops range report X-721-67-103, is over-predicted by
roughly $+24\%$ to $+36\%$, and the Nike--Cajun sounding flight (NACA RM~L57D26)
is over-predicted by $+16.6\%$. The original blunt and secant-ogive Arcas flights
(DTIC AD-235341) and the HEROS~3 amateur flight (Kobald, 2018) are instead
severely under-predicted, by $-29\%$ to $-69\%$. Two cases are outright
simulation failures rather than
honest mispredictions: Nike--Apache 14.210~GI returns a degenerate zero apogee
(a launch-state geometry/recovery configuration error), and the Nike--Cajun
Hurricane flight terminates with an integration error before apogee; both are
listed in Table~\ref{tab:exploratory} for completeness but carry no physical
interpretation.

\begin{table*}[t]
\centering
\caption{Full exploratory high-Mach sounding-rocket set (reported in its entirety;
none of these flights enters any headline statistic of Sec.~8). Signed apogee
error is (predicted~$-$~measured)/measured. Rows marked \textsuperscript{$\dagger$}
are simulation failures, not physical predictions.}
\label{tab:exploratory}
\begin{tabular}{l c S[table-format=1.3] S[table-format=+3.2] c}
\toprule
Flight & Source & {Peak $M$} & {Error (\%)} & Within $\pm10\%$ \\
\midrule
\multicolumn{5}{l}{\textit{Within $\pm10\%$ (well-characterized vehicles)}} \\
Black Brant~V~VB (AAF-VB-32)        & \cite{dtic_ad0733141}  & 7.224 & -6.97 & yes \\
Nike--Deacon no.~1 (Wallops 1955)   & \cite{heitkotter1956}  & 4.956 & -1.06 & yes \\
Nike--Deacon no.~2 (Wallops 1955)   & \cite{heitkotter1956}  & 5.079 & -0.89 & yes \\
\midrule
\multicolumn{5}{l}{\textit{Over-predicted (Nike--Apache 1965 set; Nike--Cajun)}} \\
Nike--Apache 14.108~GI              & \cite{nasa_x721_66_568} & 6.498 & +31.47 & no \\
Nike--Apache 14.75~GR (Wallops 1965)& X-721-67-103            & 6.747 & +24.49 & no \\
Nike--Apache 14.78~UA (WSMR 1965)   & X-721-67-103            & 5.955 & +23.87 & no \\
Nike--Apache 14.133~NA (Wallops 1965)& X-721-67-103           & 6.335 & +33.82 & no \\
Nike--Apache 14.213~UI (Wallops 1965)& X-721-67-103           & 6.426 & +35.85 & no \\
Nike--Apache 14.214~UI (Wallops 1965)& X-721-67-103           & 6.416 & +35.32 & no \\
Nike--Apache 14.242~UE (Wallops 1965)& X-721-67-103           & 7.047 & +32.68 & no \\
Nike--Apache 14.244~UI (Wallops 1965)& X-721-67-103           & 6.939 & +26.88 & no \\
Nike--Apache 14.247~UI (Wallops 1965)& X-721-67-103           & 7.047 & +29.63 & no \\
Nike--Cajun UM (Wallops 1956--57)   & NACA RM L57D26          & 6.249 & +16.56 & no \\
\midrule
\multicolumn{5}{l}{\textit{Under-predicted (Arcas family; HEROS~3)}} \\
Arcas flight 2 (blunt)              & DTIC AD-235341          & 2.294 & -29.18 & no \\
Arcas flight 3 (blunt)              & DTIC AD-235341          & 2.294 & -38.76 & no \\
Arcas flight 4 (secant ogive)       & DTIC AD-235341          & 2.535 & -69.12 & no \\
Arcas flight 5 (secant ogive)       & DTIC AD-235341          & 2.535 & -67.93 & no \\
HEROS~3 (Esrange 2016)              & Kobald 2018             & 1.886 & -63.38 & no \\
\midrule
\multicolumn{5}{l}{\textit{Simulation failures (no physical prediction)}} \\
Nike--Apache 14.210~GI\textsuperscript{$\dagger$} & X-721-67-103 & {---} & {zero apogee} & no \\
Nike--Cajun Hurricane\textsuperscript{$\dagger$}  & NACA RM L57D26 & {---} & {integration error} & no \\
\bottomrule
\end{tabular}
\end{table*}

\subsection{Root-Cause Assessment}

The dominant source of error in this set is \emph{vehicle reconstruction}, not a
single defect in the aerodynamic models. The headline corpus of Sec.~8 pairs each
flight with a manufacturer or contest-documented thrust curve and a complete
airframe definition; the historical sounding rockets here do not. Their thrust
histories were synthesized from total-impulse and burn-time figures, and their
external geometries were reconstructed from partial drawings and performance
handbooks. The two opposite-signed failure modes are consistent with this
diagnosis rather than with a common modelling bias. The Arcas and HEROS~3
under-predictions, which occur at modest peak Mach numbers ($1.9$--$2.5$), are the
signature of under-stated propulsive energy in the reconstructed motors---an
178,000~ft documented Arcas apogee cannot be matched by any drag model when the
synthesized impulse is too low, and the error there is propulsive bookkeeping, not
high-Mach aerodynamics.

The Nike--Apache over-predictions are more interesting because they are
\emph{systematic and same-signed} across nine independent flights ($+24\%$ to
$+36\%$), all peaking above Mach~$5.9$. A purely random reconstruction error would
not bias an entire vehicle family in one direction. This pattern instead points to
an aerodynamic or flight-dynamics mismatch specific to this configuration above
Mach~$5$---most plausibly an under-predicted coast-phase drag coefficient (too
little aerodynamic deceleration during the long high-altitude coast to apogee),
together with the simplified thrust-misalignment and dispersion assumptions of a
nominal vertical trajectory. Identifying which of these dominates would require the
detailed range and trajectory data that motivate treating the Nike--Apache result
as an open question rather than a validated outcome.

\subsection{Status and Future Work}

We draw no validation claim from this set. The honest reading is that the
method's supersonic models continue to run, and can be accurate to within about
$7\%$ at Mach~7, but only on the small number of flights for which the propulsion
and geometry are independently documented to the standard of the headline corpus.
On poorly documented historical vehicles, reconstruction uncertainty swamps any
model signal. Establishing genuine hypersonic flight validation therefore
requires either flights with experimentally measured thrust and mass-property
histories at these Mach numbers, or a dedicated reconstruction-uncertainty
quantification that propagates the documented tolerances on impulse, mass, and
geometry into the predicted apogee. Both are left to future work; until then the
paper's validated envelope remains the supersonic regime to Mach~$4.33$
established in Sec.~8, and the high-Mach reach demonstrated here is reported as
exploratory only.
